% ===========================================================================
\section{Introduction}
% Source material: README.md, docs/architecture.md
% ===========================================================================

Institutional elections usually run on one operator's database. Voters have to
trust that operator with the voter roll, the ballot box, the count and the log,
and nobody outside can check any of them.

\projectname{} is a voting application that takes the operator out of the ballot
box and the count. Smart contracts on a public test blockchain hold the roll,
the schedule and the count. The Election Authority cannot cast or change a
ballot, change the roll after registration closes, or stop the count.

At registration, a voter creates a secret on their own device and hands the
registrar only a commitment, a hash of that secret. When registration closes,
the registrar puts the list of commitments on IPFS and writes two things
on-chain: the Merkle root of that list and the list's IPFS address. To vote,
the voter's browser downloads the list and builds a zero-knowledge proof that
the voter owns one of the commitments, without revealing which one. The
registrar is not contacted at vote time.

Each ballot carries the proof, the chosen candidate, and a nullifier,
\texttt{hash(identitySecret, electionId)}. The nullifier is the same every time a
voter tries to vote in the same election, so the contract rejects a second
ballot from the same voter. The contract also stores the number of registered
voters and refuses any ballot beyond it. The tally is contract state and
changes only when a ballot is accepted. Anyone can check every proof and
recount the totals without trusting the Election Authority, the relayer, or us.

We chose voting because it has the widest gap between what the operator claims
and what anyone can check. A wrong stock level is caught at the next physical
count; a wrong election result has no such later check. In an election the
participants are opponents, so the record must be one that no participant
controls, that nobody can edit, and that anyone can check. An audit does not
fix this, because the audit trail is kept by the party being audited. A public
blockchain does: the rules are enforced by the network, and anyone can re-run
them.

The design does not solve every problem in electronic voting. Its limits (roll
correctness, coercion, the trusted setup, and privacy from the relayer) are
listed once, with reasons, in Section~\ref{sec:trust}.

Section~2 looks at the problem; Section~3 gives the requirements and the trust
model; Section~4 the design and technology choices; Section~5 the evaluation
parameters; Section~6 the plan and risks. Nothing is built yet: every number in
Section~5 is a target for Evaluation~2, not a measurement.
